\documentclass[aspectratio=169,11pt]{beamer}

\usetheme{Madrid}
\usecolortheme{whale}
\setbeamertemplate{navigation symbols}{}
\setbeamertemplate{footline}[frame number]

\usepackage[utf8]{inputenc}
\usepackage[T1]{fontenc}
\usepackage{graphicx}
\usepackage{booktabs}
\usepackage{amsmath}
\usepackage{hyperref}
\usepackage{tikz}
\usetikzlibrary{shapes,arrows,positioning,calc,fit,backgrounds}
\usepackage{xcolor}
\usepackage{array}
\usepackage{listings}
\usepackage{fontawesome5}

\lstset{
    basicstyle=\ttfamily\scriptsize,
    breaklines=true,
    frame=single,
    backgroundcolor=\color{gray!10},
    keywordstyle=\color{blue!70},
    commentstyle=\color{green!50!black},
    stringstyle=\color{red!60},
    showstringspaces=false,
    tabsize=2
}

\definecolor{stepbg}{RGB}{232,245,233}
\definecolor{stepborder}{RGB}{76,175,80}
\definecolor{aibg}{RGB}{243,229,245}
\definecolor{aiborder}{RGB}{156,39,176}
\definecolor{verifybg}{RGB}{227,242,253}
\definecolor{verifyborder}{RGB}{33,150,243}
\definecolor{hwbg}{RGB}{255,243,224}
\definecolor{hwborder}{RGB}{255,152,0}
\definecolor{tryitbg}{RGB}{255,253,231}
\definecolor{tryitborder}{RGB}{198,167,0}
\definecolor{driftbg}{RGB}{255,235,238}
\definecolor{driftborder}{RGB}{211,47,47}

\title[EECE5155 -- MLOps Seminar]{MLOps for IoT: Hands-On Seminar}
\subtitle{Real Sensor Data $\rightarrow$ Train $\rightarrow$ Deploy $\rightarrow$ Drift $\rightarrow$ Retrain}
\author{Eduardo Baena}
\date{Spring 2026}
\institute{Northeastern University\\
\texttt{e.baena@northeastern.edu}}

\begin{document}

%--- Title ---
\begin{frame}
    \titlepage
    \vspace{-0.8cm}
    \begin{center}
        \textbf{EECE5155: Wireless Sensor Networks and the Internet of Things}\\[4pt]
        {\small \faIcon{github}\ \texttt{\href{https://github.com/ebaenamar/ECE5155\_practical\_MLOps}{github.com/ebaenamar/ECE5155\_practical\_MLOps}}}
    \end{center}
\end{frame}

%==============================================================================
\section{The Real-World Problem}
%==============================================================================

\begin{frame}{The IoT Scenario: Activity Recognition on a Wearable}
    \begin{columns}
        \column{0.55\textwidth}
        \textbf{The device:} smartphone worn on the waist\\
        \textbf{The sensors:} 3-axis accelerometer + gyroscope at 50 Hz\\
        \textbf{The goal:} automatically detect what the user is doing

        \vspace{0.3cm}
        \textbf{Six activity classes:}
        \begin{enumerate}
            \item WALKING
            \item WALKING\_UPSTAIRS
            \item WALKING\_DOWNSTAIRS
            \item SITTING
            \item STANDING
            \item LAYING
        \end{enumerate}

        \column{0.42\textwidth}
        \begin{exampleblock}{Why this matters in IoT}
            Wearable health monitors, fall detection for elderly care, physical therapy tracking, fitness wearables. The same pipeline runs on Fitbit, Apple Watch, hospital wristbands.
        \end{exampleblock}

        \vspace{0.2cm}
        \begin{block}{The dataset: UCI HAR}
            30 volunteers, ages 19--48, real smartphone recordings, video-labeled.
            10,299 samples. Published 2013, cited 5,000+ times.\\
            {\footnotesize \texttt{archive.ics.uci.edu/dataset/240}}
        \end{block}
    \end{columns}
\end{frame}

\begin{frame}{What Is Training Data? --- Concretely}
    \begin{columns}
        \column{0.52\textwidth}
        \textbf{How this dataset was collected:}
        \begin{enumerate}
            \item 30 volunteers wore a phone on their waist
            \item Each performed 6 activities while \textbf{video-recorded}
            \item Researchers watched the video and wrote down which activity happened at every moment
            \item Those timestamps became the \textbf{labels}
        \end{enumerate}

        \vspace{0.2cm}
        \textbf{One labeled example =}\\
        \hspace{0.3cm} 2.56 s of sensor readings (128 samples)\\
        \hspace{0.3cm} + one label: ``WALKING''

        \column{0.45\textwidth}
        \begin{alertblock}{Key engineering insight}
            The model never sees raw signal. From each 2.56-second window, \textbf{561 statistical features} are extracted: mean, std, FFT coefficients, correlations between axes. This is \textbf{feature engineering} --- the same step in your sensor pipeline before running ML.
        \end{alertblock}

        \vspace{0.2cm}
        \begin{exampleblock}{The labeling cost}
            30 people $\times$ hours of recording $\times$ manual video review. This is why labeled IoT data is expensive. In production, a domain expert (nurse, engineer) labels every event.
        \end{exampleblock}
    \end{columns}
\end{frame}

\begin{frame}{Seminar Timeline}
    \begin{columns}
        \column{0.55\textwidth}
        \textbf{What you will build today:}
        \begin{enumerate}
            \item \texttt{explore.py} --- visualize the data
            \item \texttt{train.py} --- train Random Forest, learn the split trap
            \item \texttt{compare\_models.py} --- kNN vs RF vs SVM: accuracy, speed, size
            \item \texttt{drift\_sim.py} --- deploy to elderly patients, watch accuracy drop
            \item \texttt{retrain.py} --- collect new data, recover accuracy
        \end{enumerate}

        \column{0.42\textwidth}
        \begin{alertblock}{Timeline}
            \textbf{0--5 min:} Setup\\
            \textbf{5--20 min:} Explore + train\\
            \textbf{20--35 min:} Compare 3 algorithms\\
            \textbf{35--50 min:} Drift + retrain\\
            \textbf{50--60 min:} Take-home briefing
        \end{alertblock}

        \vspace{0.2cm}
        \begin{exampleblock}{Grading}
            Working code only: \textbf{50\%}\\
            Broken code + debug log: \textbf{80\%}\\
            Working + reasoning + AI audit: \textbf{100\%}
        \end{exampleblock}
    \end{columns}
\end{frame}

%==============================================================================
\section{Step 0: Setup}
%==============================================================================

\begin{frame}[fragile]{Step 0a: Install Dependencies (Before Class)}
    \fcolorbox{stepborder}{stepbg}{\begin{minipage}{0.93\textwidth}
    \textbf{\faIcon{check-circle} Do this before the seminar. Takes 5 minutes.}
    \end{minipage}}

    \vspace{0.15cm}
    \textbf{1. Check you have Python 3.9+:}
\begin{lstlisting}[language=bash]
python --version          # Windows / Linux
python3 --version         # macOS
\end{lstlisting}

    \vspace{0.1cm}
    \textbf{2. Install dependencies (same command on all OS):}
\begin{lstlisting}[language=bash]
pip install scikit-learn pandas numpy matplotlib seaborn scipy
# macOS/Linux: use pip3 if pip is not found
# Windows: open Anaconda Prompt or PowerShell as Administrator
\end{lstlisting}

    \vspace{0.1cm}
    \textbf{3. Verify install:}
\begin{lstlisting}[language=Python]
import sklearn, pandas, numpy, matplotlib, seaborn, scipy
print("All packages OK")
\end{lstlisting}

    \fcolorbox{verifyborder}{verifybg}{\begin{minipage}{0.93\textwidth}
    \textbf{\faIcon{check} If you see ``All packages OK'' with no errors, you are ready for step 0b.}
    \end{minipage}}
\end{frame}

\begin{frame}[fragile,shrink=15]{Step 0b: Download the Dataset (Run Once)}
    \fcolorbox{stepborder}{stepbg}{\begin{minipage}{0.93\textwidth}
    \textbf{\faIcon{download} UCI HAR dataset: $\sim$60 MB, saves 4 CSV files. Run once before class.}
    \end{minipage}}

    \vspace{0.05cm}
    \textbf{Download and save as CSV files:}
\begin{lstlisting}[language=Python,basicstyle=\ttfamily\tiny]
import urllib.request, zipfile, io, pandas as pd

URL = ("https://archive.ics.uci.edu/static/public/240/"
       "human+activity+recognition+using+smartphones.zip")
req = urllib.request.Request(URL, headers={"User-Agent": "Mozilla/5.0"})
with urllib.request.urlopen(req, timeout=120) as r:
    outer = zipfile.ZipFile(io.BytesIO(r.read()))
inner = zipfile.ZipFile(io.BytesIO(outer.read("UCI HAR Dataset.zip")))

def read(path):
    with inner.open(path) as f: return f.read().decode()

features = [l.split(" ",1)[1].strip()
            for l in read("UCI HAR Dataset/features.txt").strip().split("\n")]
labels   = {int(l.split()[0]): l.split()[1]
            for l in read("UCI HAR Dataset/activity_labels.txt").strip().split("\n")}

def load(split):
    X = pd.read_csv(io.StringIO(read(f"UCI HAR Dataset/{split}/X_{split}.txt")),
                    sep=r"\s+", header=None, names=features)
    y = pd.read_csv(io.StringIO(read(f"UCI HAR Dataset/{split}/y_{split}.txt")),
                    header=None)[0].map(labels).rename("Activity")
    return X, y

X_train, y_train = load("train")
X_test,  y_test  = load("test")
X_train.to_csv("har_X_train.csv", index=False)
X_test.to_csv( "har_X_test.csv",  index=False)
y_train.to_csv("har_y_train.csv", index=False)
y_test.to_csv( "har_y_test.csv",  index=False)
print(f"Train: {X_train.shape}   Test: {X_test.shape}")
print(y_train.value_counts())
\end{lstlisting}

    \vspace{0.05cm}
    \fcolorbox{verifyborder}{verifybg}{\begin{minipage}{0.93\textwidth}
    \textbf{\faIcon{check} Expected: \texttt{Train: (7352, 561) \ \ Test: (2947, 561)} and 6 activity classes listed. You are ready.}
    \end{minipage}}
\end{frame}

%==============================================================================
\section{Part 1: Explore and Train (5--20 min)}
%==============================================================================

\begin{frame}[shrink=5]{What Are We Actually Doing?}
    \begin{columns}
        \column{0.52\textwidth}
        \textbf{The core idea of supervised learning:}
        \begin{enumerate}
            \item Collect sensor data where you \textbf{already know the answer}\\(e.g.\ a video confirms ``this person is walking'')
            \item Show the algorithm many labeled examples\\$\rightarrow$ this is called \textbf{training}
            \item Show it \textbf{new examples it has never seen}\\$\rightarrow$ measure how often it gets the right answer\\$\rightarrow$ this is called \textbf{testing}
        \end{enumerate}

        \vspace{0.2cm}
        \textbf{Why split into train and test?}\\Because if you test on the same data you trained on, you are asking ``do you remember?'', not ``do you understand?''. A student who memorized the exam answers scores 100\% but learned nothing.

        \column{0.45\textwidth}
        \begin{exampleblock}{In numbers for this dataset}
            \textbf{7352 examples} $\rightarrow$ the model learns from these\\[4pt]
            \textbf{2947 examples} $\rightarrow$ locked away during training, used only to measure accuracy\\[4pt]
            The split was made by \textbf{person}: subjects 1--21 train, subjects 22--30 test. Never mixed.
        \end{exampleblock}

        \vspace{0.2cm}
        \begin{alertblock}{One number to understand}
            \textbf{Accuracy = correct predictions / total predictions}\\If the model predicts 2700 out of 2947 correctly: accuracy = 2700/2947 = 0.916
        \end{alertblock}
    \end{columns}
\end{frame}

\begin{frame}[fragile,shrink=5]{Step 1: Visualize the Data}
    \fcolorbox{stepborder}{stepbg}{\begin{minipage}{0.93\textwidth}
    \textbf{\faIcon{search} Before training anything, always look at the data. Create \texttt{explore.py}.}
    \end{minipage}}

    \vspace{0.05cm}
    {\footnotesize \texttt{read\_csv} loads the table. \texttt{X} = 561 numbers/sample. \texttt{y} = label. \texttt{value\_counts()} counts samples/class. \texttt{scatter()} plots 2 features coloured by activity.}

\begin{lstlisting}[language=Python,basicstyle=\ttfamily\tiny]
import pandas as pd, matplotlib.pyplot as plt, seaborn as sns

X_train = pd.read_csv("har_X_train.csv")
X_test  = pd.read_csv("har_X_test.csv")
y_train = pd.read_csv("har_y_train.csv").squeeze()
y_test  = pd.read_csv("har_y_test.csv").squeeze()

# 1. Class balance (train set)
y_train.value_counts().plot(kind="bar", title="Samples per activity", color="steelblue")
plt.tight_layout(); plt.savefig("class_balance.png"); plt.close()

# 2. Scatter: two features that strongly separate activities
fig, ax = plt.subplots(figsize=(9, 4))
palette = sns.color_palette("tab10", 6)
X_all = pd.concat([X_train, X_test], ignore_index=True)
y_all = pd.concat([y_train, y_test], ignore_index=True)
for i, act in enumerate(y_all.unique()):
    mask = (y_all == act)
    ax.scatter(X_all.loc[mask, "tBodyAccMag-mean()"],
               X_all.loc[mask, "tGravityAcc-mean()-X"],
               label=act, alpha=0.3, s=10, color=palette[i])
ax.set_xlabel("Body Acc Magnitude (mean)")
ax.set_ylabel("Gravity Acc X (mean)")
ax.legend(fontsize=7, ncol=2)
plt.tight_layout(); plt.savefig("feature_scatter.png"); plt.close()
print("Plots saved.")
\end{lstlisting}

    \fcolorbox{verifyborder}{verifybg}{\begin{minipage}{0.93\textwidth}
    \textbf{\faIcon{check} Scatter: LAYING separate (lying flat). WALKING separate (high energy). SITTING/STANDING \emph{overlap} --- gravity vectors nearly identical for both postures.}
    \end{minipage}}
\end{frame}

\begin{frame}[shrink=10]{What Is a Random Forest?}
    \begin{columns}
        \column{0.52\textwidth}
        \textbf{One decision tree = a chain of yes/no questions:}
        \begin{enumerate}
            \item Is Body Acc Magnitude $>$ -0.5?
            \item If yes: is Gravity X $>$ 0.3? $\rightarrow$ WALKING
            \item If no: is Gyro Z mean $<$ 0.1? $\rightarrow$ LAYING
            \item \ldots and so on for each of the 561 features
        \end{enumerate}

        \vspace{0.2cm}
        \textbf{Random Forest = 100 trees, each trained on a random subset of the data.} Final answer = majority vote of 100 trees. More trees $\rightarrow$ more robust, less likely to memorize noise.

        \column{0.45\textwidth}
        \begin{exampleblock}{What \texttt{rf.fit(X\_train, y\_train)} does}
            The algorithm looks at all 7352 training examples and builds 100 decision trees. Each tree learns which combinations of the 561 features distinguish the 6 activities. This step can take 10--30 seconds. \textbf{You only do this once.}
        \end{exampleblock}

        \vspace{0.2cm}
        \begin{exampleblock}{What \texttt{rf.predict(X\_test)} does}
            For each of the 2947 test examples, the 100 trees each vote for one activity. The majority vote wins. This is \textbf{inference} --- it runs in milliseconds.
        \end{exampleblock}
    \end{columns}
\end{frame}

\begin{frame}[shrink=8]{How Do We Measure Whether a Model Works?}
    \begin{columns}
        \column{0.50\textwidth}
        \textbf{Accuracy} --- the obvious metric:
        \[\text{accuracy} = \frac{\text{correct predictions}}{\text{total predictions}}\]
        \textbf{Problem:} a fall-detection model that \emph{always} predicts ``no fall'' achieves 99\% accuracy --- and misses every real fall. In IoT, different errors have very different costs.

        \vspace{0.2cm}
        \begin{block}{Two metrics that matter more}
            \textbf{Precision:} of all the times the model fired the alarm, how many were real?\\
            $\rightarrow$ high precision = few false alarms\\[4pt]
            \textbf{Recall:} of all real events, how many did we catch?\\
            $\rightarrow$ high recall = few missed events\\[4pt]
            \textbf{F1:} harmonic mean of both --- one summary number per class
        \end{block}

        \column{0.47\textwidth}
        \textbf{Confusion matrix} --- the most informative output:

        \vspace{0.1cm}
        {\small Rows = true label. Columns = predicted label. Diagonal = correct. Off-diagonal = errors and exactly what got confused with what.}

        \vspace{0.2cm}
        \begin{alertblock}{IoT framing}
            {\small Which failure modes are you shipping?\\
            A fall-detector confusing WALKING with FALLING $\rightarrow$ constant alarms in the corridor.\\
            A maintenance system confusing ``early fault'' with ``normal'' $\rightarrow$ misses the event it was built to catch.}
        \end{alertblock}

        \vspace{0.15cm}
        \begin{exampleblock}{Workflow}
            {\small 1. Accuracy --- ballpark check\\
            2. Per-class recall --- critical classes detected?\\
            3. Confusion matrix --- which errors?\\
            4. Are those errors acceptable for your use case?}
        \end{exampleblock}
    \end{columns}
\end{frame}

\begin{frame}[fragile,shrink=5]{Step 2: Train Your First Model}
    \fcolorbox{stepborder}{stepbg}{\begin{minipage}{0.93\textwidth}
    \textbf{\faIcon{python} Create \texttt{train.py} --- split by subject, train Random Forest.}
    \end{minipage}}

\begin{lstlisting}[language=Python]
import pandas as pd
from sklearn.ensemble import RandomForestClassifier
from sklearn.metrics import accuracy_score, classification_report
import pickle

X_train = pd.read_csv("har_X_train.csv")
X_test  = pd.read_csv("har_X_test.csv")
y_train = pd.read_csv("har_y_train.csv").squeeze()
y_test  = pd.read_csv("har_y_test.csv").squeeze()
print(f"Train: {len(X_train)} | Test: {len(X_test)}")

rf = RandomForestClassifier(n_estimators=100, random_state=42, n_jobs=-1)
rf.fit(X_train, y_train)

acc = accuracy_score(y_test, rf.predict(X_test))
print(f"\nRandom Forest accuracy: {acc:.3f}")
print(classification_report(y_test, rf.predict(X_test)))
pickle.dump(rf, open("rf_model.pkl", "wb"))
\end{lstlisting}

    \fcolorbox{tryitborder}{tryitbg}{\begin{minipage}{0.93\textwidth}
    \textbf{\faIcon{hand-point-right} After running: which two activities are most confused? Check the classification report. Why does that make physical sense?}
    \end{minipage}}
\end{frame}

\begin{frame}[shrink=5]{Why Does the Split Strategy Matter?}
    \begin{columns}
        \column{0.52\textwidth}
        \textbf{Imagine this exam scenario:}\\[4pt]
        A professor gives students the exam questions one week early. Students memorize the answers. On exam day, everyone scores 100\%. Does that mean they learned the material?\\[8pt]
        \textbf{The same thing happens in ML} when you mix data from the same person in both train and test sets:
        \begin{itemize}
            \item The model sees subject 5's walking pattern during training
            \item It is tested on \textbf{another window from subject 5}
            \item It ``recognizes'' the person, not the activity
            \item Accuracy looks great --- but it fails on any new user
        \end{itemize}

        \column{0.45\textwidth}
        \begin{alertblock}{The correct question to ask}
            ``If I deploy this to a new person the model has \textbf{never seen}, how well does it work?''\\[4pt]
            This is why the UCI HAR dataset splits by \textbf{subject}: the 9 test subjects were never in the training set.
        \end{alertblock}

        \vspace{0.2cm}
        \begin{exampleblock}{What you will observe}
            Subject split: \textbf{0.929} $\leftarrow$ real accuracy\\[2pt]
            Random shuffle: \textbf{0.978} $\leftarrow$ fake, do not trust\\[4pt]
            A 5 pp gap that would mislead any engineering decision.
        \end{exampleblock}
    \end{columns}
\end{frame}

\begin{frame}[fragile,shrink=5]{Step 3: The Split Trap}
    \fcolorbox{driftborder}{driftbg}{\begin{minipage}{0.93\textwidth}
    \textbf{\faIcon{exclamation-triangle} A critical evaluation mistake. Add this to \texttt{train.py} and run it.}
    \end{minipage}}

    \vspace{0.1cm}
\begin{lstlisting}[language=Python]
from sklearn.model_selection import train_test_split

# Wrong split: random 80/20 shuffle across all subjects
X_all = pd.concat([X_train, X_test], ignore_index=True)
y_all = pd.concat([y_train, y_test], ignore_index=True)
X_tr2, X_te2, y_tr2, y_te2 = train_test_split(
    X_all, y_all, test_size=0.2, random_state=42)
rf2 = RandomForestClassifier(n_estimators=100, random_state=42, n_jobs=-1)
rf2.fit(X_tr2, y_tr2)
acc_wrong = accuracy_score(y_te2, rf2.predict(X_te2))
print(f"Subject split:  {acc:.3f}   <-- real deployment accuracy")
print(f"Random split:   {acc_wrong:.3f}   <-- inflated, do not trust!")
\end{lstlisting}

    \vspace{0.1cm}
    \textbf{Expected output:}
\begin{lstlisting}
Subject split:  0.929   <-- real deployment accuracy
Random split:   0.977   <-- inflated, do not trust!
\end{lstlisting}

    \fcolorbox{verifyborder}{verifybg}{\begin{minipage}{0.93\textwidth}
    \textbf{Why the 5\% gap?} Random split leaks windows from the same subject into both train and test. The model memorizes per-person gait patterns. When deployed to \textbf{new users} it cannot see, the real accuracy is 0.93, not 0.98. \textbf{Your evaluation must match your deployment scenario.}
    \end{minipage}}
\end{frame}

%==============================================================================
\section{Part 2: Compare ML Methods (20--35 min)}
%==============================================================================

\begin{frame}[shrink=10]{Part 2: Choosing the Right Algorithm for IoT}
    \begin{columns}
        \column{0.55\textwidth}
        \textbf{Three candidates on this problem:}

        \vspace{0.2cm}
        \begin{tabular}{lp{4.2cm}}
            \toprule
            \textbf{Method} & \textbf{Core idea} \\
            \midrule
            k-NN & find the $k$ nearest training examples; majority vote \\
            Random Forest & 100 decision trees; majority vote \\
            SVM (linear) & hyperplane maximizing class margins \\
            \bottomrule
        \end{tabular}

        \vspace{0.3cm}
        \textbf{The IoT constraint nobody mentions:}\\
        A microcontroller has 256 kB RAM. k-NN stores the entire training set: $7352 \times 561$ floats = \textbf{16 MB}. That does not fit on an MCU. Accuracy alone is not enough.

        \column{0.42\textwidth}
        \begin{alertblock}{What you are really choosing}
            For each algorithm, you care about:\\[4pt]
            \textbf{Accuracy} --- does it work?\\
            \textbf{Inference time} --- does it run at 50 Hz?\\
            \textbf{Model size} --- does it fit in flash?\\
            \textbf{Training data needed} --- how many labels?\\[4pt]
            Today you will measure all four.
        \end{alertblock}
    \end{columns}
\end{frame}

\begin{frame}[fragile]{Why Do We Need to ``Scale'' the Data?}
    \begin{columns}
        \column{0.52\textwidth}
        \textbf{The problem with raw features:}\\[4pt]
        The 561 features have very different numerical ranges:
        \begin{itemize}
            \item \texttt{tBodyAcc-mean()-X} $\approx$ $-$0.3 to $+$0.3
            \item \texttt{tGravityAcc-energy()-Z} $\approx$ $-$1.0 to $+$1.0
            \item \texttt{fBodyAcc-maxInds-X} $\approx$ $-$128 to $0$
        \end{itemize}
        \vspace{0.1cm}
        A model that compares distances (k-NN) or maximizes margins (SVM) is dominated by the feature with the largest numbers. That is an accident of units, not physics.\\[4pt]
        \textbf{StandardScaler}: subtract mean, divide by std $\rightarrow$ every feature has mean=0, std=1.

        \column{0.45\textwidth}
        \textbf{Important rule --- fit only on train:}
\begin{lstlisting}[language=Python]
scaler = StandardScaler()
scaler.fit(X_train)   # learn mean/std
Xtr = scaler.transform(X_train)
Xte = scaler.transform(X_test)
# Never fit on X_test: that leaks
# information about the future.
\end{lstlisting}
        \vspace{0.15cm}
        \begin{block}{Random Forest does not need scaling}
            Trees split on individual feature thresholds --- absolute scale does not matter. k-NN and SVM do need it.
        \end{block}
    \end{columns}
\end{frame}

\begin{frame}[fragile,shrink=10]{Step 4: Compare Three Classifiers}
    \fcolorbox{stepborder}{stepbg}{\begin{minipage}{0.93\textwidth}
    \textbf{\faIcon{python} Create \texttt{compare\_models.py}}
    \end{minipage}}

\begin{lstlisting}[language=Python,basicstyle=\ttfamily\tiny]
import time, os, pickle, tempfile, pandas as pd
from sklearn.ensemble import RandomForestClassifier
from sklearn.neighbors import KNeighborsClassifier
from sklearn.svm import LinearSVC
from sklearn.preprocessing import StandardScaler
from sklearn.metrics import accuracy_score

X_train = pd.read_csv("har_X_train.csv")
X_test  = pd.read_csv("har_X_test.csv")
y_train = pd.read_csv("har_y_train.csv").squeeze()
y_test  = pd.read_csv("har_y_test.csv").squeeze()

scaler  = StandardScaler().fit(X_train)
Xtr_sc  = scaler.transform(X_train)
Xte_sc  = scaler.transform(X_test)

models = {
    "k-NN (k=5)":    (KNeighborsClassifier(n_neighbors=5, n_jobs=-1), True),
    "Random Forest": (RandomForestClassifier(n_estimators=100, random_state=42, n_jobs=-1), False),
    "SVM (linear)":  (LinearSVC(max_iter=2000, random_state=42), True),
}
for name, (clf, use_scale) in models.items():
    Xtr = Xtr_sc if use_scale else X_train
    Xte = Xte_sc if use_scale else X_test
    t0 = time.time();  clf.fit(Xtr, y_train);  train_t = time.time()-t0
    t0 = time.time();  preds = clf.predict(Xte)
    infer_ms = (time.time()-t0) / len(Xte) * 1000
    acc = accuracy_score(y_test, preds)
    with tempfile.NamedTemporaryFile(delete=False) as f:
        pickle.dump(clf, f);  size_kb = os.path.getsize(f.name)/1024
        os.unlink(f.name)
    print(f"{name:<22} acc={acc:.3f}  train={train_t:.1f}s"
          f"  infer={infer_ms:.3f}ms  size={size_kb:.0f}kB")
\end{lstlisting}
\end{frame}

\begin{frame}[fragile,shrink=10]{Step 5: Read the Comparison and Decide}
    \fcolorbox{stepborder}{stepbg}{\begin{minipage}{0.93\textwidth}
    \textbf{\faIcon{play} Run: \texttt{python compare\_models.py}}
    \end{minipage}}

    \vspace{0.2cm}
    \textbf{Expected output (your times will vary):}
\begin{lstlisting}
k-NN (k=5)             acc=0.884  train=0.1s  infer=~8ms  size=16384kB
Random Forest          acc=0.929  train=~12s  infer=~0.05ms  size=5824kB
SVM (linear)           acc=0.963  train=~4s   infer=~0.01ms  size=27kB
\end{lstlisting}

    \vspace{0.15cm}
    \begin{center}
    \footnotesize
    \begin{tabular}{lcccc}
        \toprule
        \textbf{Method} & \textbf{Accuracy} & \textbf{Infer/sample} & \textbf{Model size} & \textbf{MCU deployable?} \\
        \midrule
        k-NN           & 0.884 & $\sim$8 ms & 16 MB (stores all data)  & No \\
        Random Forest  & 0.929 & $\sim$0.05 ms & 5.7 MB (100 trees) & No \\
        SVM (linear)   & \textbf{0.963} & \textbf{$\sim$0.01 ms} & \textbf{27 kB} & \textbf{Yes} \\
        \bottomrule
    \end{tabular}
    \normalsize
    \end{center}

    \fcolorbox{tryitborder}{tryitbg}{\begin{minipage}{0.93\textwidth}
    \textbf{\faIcon{hand-point-right} Three things to discuss:}\\
    1. Why is k-NN inference 600$\times$ slower than SVM?\\
    2. Why does SVM win on accuracy here, despite being the simplest model?\\
    3. For a Raspberry Pi gateway (no MCU constraint), which would you pick?
    \end{minipage}}
\end{frame}

%==============================================================================
\section{Part 3: Drift and Retrain (35--50 min)}
%==============================================================================

\begin{frame}[shrink=8]{Part 3: Concept Drift on a Wearable}
    \begin{columns}
        \column{0.55\textwidth}
        \textbf{The scenario:}\\
        Your SVM was trained on healthy adults aged 19--48. On day one of deployment to a hospital with elderly patients, accuracy is already lower --- the model was \textbf{never trained on this population}. This is a \textbf{distribution mismatch}, not a gradual drift.

        \vspace{0.2cm}
        \textbf{Why does the model underperform on a new population?}
        \begin{itemize}
            \item Elderly gait has shorter steps and lower acceleration amplitude --- the feature distribution is different from day one
            \item Phone placement varies patient to patient: waist vs chest vs pocket
            \item Post-surgery or mobility-impaired patients move differently than healthy volunteers
            \item True temporal drift also exists: MEMS sensors age, firmware updates change signal processing
        \end{itemize}

        \column{0.42\textwidth}
        \textbf{The simulation:}\\
        \begin{exampleblock}{Why scale 0.75?}
        Menz et al.\ (2003) measured 10--20\% lower gait acceleration amplitude in elderly adults vs young adults. We use 0.75 to simulate a new deployment population. 345 feature columns contain ``Acc'' and are all scaled. This is a physics-based distribution shift, not random noise.
        \end{exampleblock}

        \vspace{0.1cm}
        \begin{block}{Which activities drift most?}
            WALKING classes --- acceleration-dependent.\\
            SITTING/STANDING/LAYING --- gravity-dominated, mostly stable.
        \end{block}
    \end{columns}
\end{frame}

\begin{frame}[shrink=8]{What Are We Actually Doing in \texttt{drift\_sim.py}?}
    \begin{columns}
        \column{0.52\textwidth}
        \textbf{The simulation in plain English:}\\[4pt]
        The model was trained on data from young adults. We want to know what happens when it is deployed to elderly patients.\\[4pt]
        We do not have real elderly patient data --- so we simulate it:\\[4pt]
        \begin{enumerate}
            \item Take the test set (young adults)
            \item Multiply all \textbf{acceleration features} by 0.75\\= ``these patients move with 25\% less amplitude''
            \item Add tiny random noise (sensor variability)
            \item Run the model on this modified data
            \item Measure how much accuracy drops
        \end{enumerate}
        \vspace{0.1cm}
        \textbf{We change the data. We do not change the model.}\\That is the whole point: the model is frozen in time.

        \column{0.45\textwidth}
        \begin{exampleblock}{The loop in the code}
            We run this for scales 1.0, 0.95, 0.90 \ldots\ 0.70. Scale 1.0 = original data. Scale 0.75 = elderly gait simulation. The plot shows how accuracy degrades as the population drifts further from training data.
        \end{exampleblock}
        \vspace{0.2cm}
        \begin{alertblock}{The key result to look for}
            At scale=0.75: accuracy is \textbf{0.742} from the first prediction on this population (baseline was 0.963). WALKING recall is 0.38. The model silently underperforms --- \textbf{with no error message, no crash, no warning}.
        \end{alertblock}
    \end{columns}
\end{frame}

\begin{frame}[fragile,shrink=20]{Step 6: Simulate Drift and Observe Degradation}
    \fcolorbox{stepborder}{stepbg}{\begin{minipage}{0.93\textwidth}
    \textbf{\faIcon{python} Create \texttt{drift\_sim.py}}
    \end{minipage}}

\begin{lstlisting}[language=Python,basicstyle=\ttfamily\tiny]
import numpy as np, pandas as pd, matplotlib.pyplot as plt
from sklearn.svm import LinearSVC
from sklearn.preprocessing import StandardScaler
from sklearn.metrics import accuracy_score

X_train = pd.read_csv("har_X_train.csv")
X_test  = pd.read_csv("har_X_test.csv")
y_train = pd.read_csv("har_y_train.csv").squeeze()
y_test  = pd.read_csv("har_y_test.csv").squeeze()

scaler = StandardScaler().fit(X_train)
svm    = LinearSVC(max_iter=2000, random_state=42)
svm.fit(scaler.transform(X_train), y_train)
baseline = accuracy_score(y_test, svm.predict(scaler.transform(X_test)))
print(f"Baseline (young adults): {baseline:.3f}")

acc_cols = [c for c in X_train.columns if "Acc" in c]  # 345 acceleration features
rng = np.random.default_rng(7)
results = []
for scale in np.arange(1.0, 0.69, -0.05):
    Xd = X_test.copy()
    Xd[acc_cols] = Xd[acc_cols] * scale
    Xd = Xd + rng.normal(0, 0.03*(1.2-scale), Xd.shape)
    acc = accuracy_score(y_test, svm.predict(scaler.transform(Xd)))
    results.append({"scale": round(scale,2), "accuracy": acc})
    print(f"  scale={scale:.2f}: acc={acc:.3f}")

df_r = pd.DataFrame(results)
plt.figure(figsize=(8,4))
plt.plot(df_r["scale"], df_r["accuracy"], "o-r", label="Drifted accuracy")
plt.axhline(baseline, linestyle="--", color="blue",
            label="Baseline %.3f" % baseline)
plt.axhline(0.85, linestyle=":", color="orange", label="Alert threshold 0.85")
plt.gca().invert_xaxis()
plt.xlabel("Acc feature scale (1.0=young adults, 0.75=elderly gait)")
plt.ylabel("Accuracy");  plt.title("SVM Accuracy Under Gait Drift")
plt.legend();  plt.tight_layout()
plt.savefig("drift_plot.png");  plt.close()
print("Saved drift_plot.png")
\end{lstlisting}
\end{frame}

\begin{frame}[fragile,shrink=20]{Step 7: Retrain on Mixed Data}
    \fcolorbox{stepborder}{stepbg}{\begin{minipage}{0.93\textwidth}
    \textbf{\faIcon{python} Create \texttt{retrain.py} --- add newly labeled samples to training set.}
    \end{minipage}}

\begin{lstlisting}[language=Python,basicstyle=\ttfamily\tiny]
import numpy as np, pandas as pd
from sklearn.svm import LinearSVC
from sklearn.preprocessing import StandardScaler
from sklearn.metrics import accuracy_score

X_train = pd.read_csv("har_X_train.csv")
X_test  = pd.read_csv("har_X_test.csv")
y_train = pd.read_csv("har_y_train.csv").squeeze()
y_test  = pd.read_csv("har_y_test.csv").squeeze()

rng      = np.random.default_rng(7)
acc_cols = [c for c in X_train.columns if "Acc" in c]

# 40% of test set re-labeled under drift = simulates new patient data
X_new = X_test.copy().sample(frac=0.4, random_state=1)
y_new = y_test.loc[X_new.index]
X_new[acc_cols] = X_new[acc_cols] * 0.75
X_new += rng.normal(0, 0.03*(1.2-0.75), X_new.shape)
X_mixed = pd.concat([X_train, X_new], ignore_index=True)
y_mixed = pd.concat([y_train, y_new], ignore_index=True)
scaler_v2 = StandardScaler().fit(X_mixed)
svm_v2    = LinearSVC(max_iter=2000, random_state=42)
svm_v2.fit(scaler_v2.transform(X_mixed), y_mixed)
# Drifted test set
X_drifted = X_test.copy()
X_drifted[acc_cols] = X_drifted[acc_cols] * 0.75
X_drifted += rng.normal(0, 0.03*(1.2-0.75), X_drifted.shape)
scaler_v1 = StandardScaler().fit(X_train)
svm_v1    = LinearSVC(max_iter=2000, random_state=42)
svm_v1.fit(scaler_v1.transform(X_train), y_train)
acc_v1 = accuracy_score(y_test, svm_v1.predict(scaler_v1.transform(X_drifted)))
acc_v2 = accuracy_score(y_test, svm_v2.predict(scaler_v2.transform(X_drifted)))
acc_v2_clean = accuracy_score(y_test, svm_v2.predict(scaler_v2.transform(X_test)))
print(f"v1 on drifted (elderly): {acc_v1:.3f}")
print(f"v2 on drifted (elderly): {acc_v2:.3f}")
print(f"v2 on clean  (young):   {acc_v2_clean:.3f}")
\end{lstlisting}
\end{frame}

\begin{frame}[fragile,shrink=8]{Step 8: The MLOps Loop Completed}
    \fcolorbox{stepborder}{stepbg}{\begin{minipage}{0.93\textwidth}
    \textbf{\faIcon{play} Run: \texttt{python retrain.py}}
    \end{minipage}}

    \vspace{0.15cm}
    \textbf{Expected results:}
    \begin{center}
    \begin{tabular}{lcc}
        \toprule
        \textbf{Scenario} & \textbf{SVM v1 (young adults)} & \textbf{SVM v2 (retrained)} \\
        \midrule
        Clean test (young adults) & 0.963 & $\sim$0.977 \\
        Drifted test (elderly gait, scale=0.75) & $\sim$0.742 & $\sim$0.976 \\
        \bottomrule
    \end{tabular}
    \end{center}

    \vspace{0.1cm}
    \fcolorbox{driftborder}{driftbg}{\begin{minipage}{0.93\textwidth}
    \textbf{What you observed:} v1 degraded from 0.963 to 0.742 with no code change, no crash, no error. v2 retrained on mixed data recovers to $\sim$0.976 on the drifted population and slightly improves on clean data (0.977). Recovery is strong because the 40\% new samples come from the same subjects as the test set --- in a real deployment, you would collect data from genuinely new patients, which would give lower but still meaningful recovery.
    \end{minipage}}

    \vspace{0.1cm}
    \fcolorbox{tryitborder}{tryitbg}{\begin{minipage}{0.93\textwidth}
    \textbf{\faIcon{hand-point-right} Discuss:} v2 gains +23 pp on the drifted population with no cost on clean data. In this experiment, new training samples come from the same test subjects --- best-case recovery. Real deployments use data from new patients, so recovery will be partial. What matters: the direction is always correct. Is mixed retraining always worth the labeling cost?
    \end{minipage}}
\end{frame}

%==============================================================================
\section{Take-Home Tasks}
%==============================================================================

\begin{frame}[shrink=5]{Take-Home Tasks}
    \fcolorbox{hwborder}{hwbg}{\begin{minipage}{0.93\textwidth}
    \textbf{\faIcon{home} Due: one week from today, end of day (Canvas).}\\
    AI is \textbf{encouraged} for take-home tasks. Always run the code to verify AI output.
    \end{minipage}}

    \vspace{0.3cm}
    \begin{center}
    \begin{tabular}{clc}
        \toprule
        \textbf{\#} & \textbf{Task} & \textbf{Weight} \\
        \midrule
        1 & Confusion matrix + feature importance analysis & 25\% \\
        2 & KS-test drift detection without labels & 35\% \\
        3 & Apply the full MLOps loop to your own project & 40\% \\
        \bottomrule
    \end{tabular}
    \end{center}

    \vspace{0.3cm}
    \textbf{For each task submit:}
    \begin{enumerate}
        \item Code + output plots or tables
        \item \textbf{AI audit:} what did you prompt? what did AI get wrong?
        \item 2--3 sentences: what did you learn?
    \end{enumerate}
\end{frame}

\begin{frame}[fragile]{Task 1: Confusion Matrix and Feature Importance (25\%)}
    \textbf{Part A: Confusion matrix} for the SVM (normalized by row = per-class recall):
    \begin{itemize}
        \item Which two activities are most confused? Write one sentence with the physical reason.
    \end{itemize}

    \vspace{0.2cm}
    \textbf{Part B: Feature importance:}
\begin{lstlisting}[language=Python]
import numpy as np
coef_importance = np.abs(svm.coef_).mean(axis=0)
top15_idx = np.argsort(coef_importance)[::-1][:15]
for i in top15_idx:
    print(f"{X.columns[i]:<42}  {coef_importance[i]:.4f}")
\end{lstlisting}

    Are the top features acceleration or gyroscope? Time domain or frequency domain? What does that tell you about which sensor and which signal property is most informative for activity recognition?

    \vspace{0.1cm}
    \fcolorbox{verifyborder}{verifybg}{\begin{minipage}{0.93\textwidth}
    \textbf{Deliverable:} Confusion matrix image + top-15 feature table + answers.
    \end{minipage}}
\end{frame}

\begin{frame}[fragile,shrink=8]{Task 2: KS-Test Drift Detection Without Labels (35\%)}
    \textbf{The real problem:} accuracy-based monitoring needs ground truth labels. In a hospital deployment, confirming ``the patient was walking'' requires manual review --- days later. You need drift detection \textbf{without labels}.

    \vspace{0.2cm}
    \textbf{Write \texttt{ks\_drift.py}} that:
    \begin{enumerate}
        \item Uses the original test set as the reference distribution
        \item Uses the drifted test set (scale 0.75) as the current distribution
        \item Runs \texttt{ks\_2samp} on all 561 features
        \item Reports: how many features flagged (p $<$ 0.05)? Compute the flagging rate separately for the 345 acceleration columns and the 213 gyroscope columns
        \item Plots a bar chart of the top 20 features by KS statistic, colored red if drifted
        \item Answers: acceleration columns are flagged at a higher rate than gyroscope-only columns. Why? (Hint: gyroscope measures rotation, not linear motion amplitude.)
    \end{enumerate}

    \fcolorbox{verifyborder}{verifybg}{\begin{minipage}{0.93\textwidth}
    \textbf{Deliverable:} Script + count table + bar chart + 3-sentence analysis.
    \end{minipage}}
\end{frame}

\begin{frame}{Task 3: Apply the Loop to Your Own Project (40\%)}
    \textbf{Write a 1-page design document for your final project:}
    \begin{enumerate}
        \item \textbf{Sensor and raw signal:} what does your device measure? At what rate?
        \item \textbf{Feature extraction:} what 4--8 features per time window? How long is the window?
        \item \textbf{Training data:} what are your classes? How do you get labeled examples? What is the bottleneck?
        \item \textbf{Algorithm choice:} using today's comparison table as reference, which algorithm fits your hardware? Justify with model size and inference time.
        \item \textbf{Expected drift causes:} list 2 specific physical reasons your distribution shifts over time.
        \item \textbf{Monitoring:} accuracy-based (if labels available) or KS-based (if not)? Justify.
    \end{enumerate}

    \begin{alertblock}{This feeds directly into your final project report}
        The design you write here becomes the MLOps section of your report.
    \end{alertblock}
\end{frame}

%==============================================================================
\section{Reference}
%==============================================================================

\begin{frame}[shrink=8]{Quick Reference Card}
    \footnotesize
    \begin{columns}
        \column{0.48\textwidth}
        \textbf{Dataset:}\\
        UCI HAR --- \texttt{archive.ics.uci.edu/dataset/240}\\
        10,299 samples $\cdot$ 561 features $\cdot$ 6 classes

        \vspace{0.15cm}
        \textbf{Files produced today:}
        \begin{tabular}{ll}
            \texttt{har\_X\_train.csv} & train features \\
            \texttt{har\_X\_test.csv}  & test features \\
            \texttt{har\_y\_train.csv} & train labels \\
            \texttt{har\_y\_test.csv}  & test labels \\
            \texttt{rf\_model.pkl}     & trained RF \\
            \texttt{class\_balance.png} & bar chart \\
            \texttt{feature\_scatter.png} & 2D scatter \\
            \texttt{drift\_plot.png}   & drift curve \\
        \end{tabular}

        \column{0.48\textwidth}
        \textbf{Key concepts from today:}
        \begin{tabular}{ll}
            Training data & labeled sensor windows \\
            Feature engineering & 561 stats from raw signal \\
            Subject split & evaluate on new users \\
            kNN vs RF vs SVM & accuracy + size + speed \\
            Concept drift & gait scale shift \\
            KS test & label-free drift detection \\
            Mixed retraining & old + new labeled data \\
        \end{tabular}

        \vspace{0.2cm}
        \textbf{References:}\\
        {[1]} Anguita et al., \textit{ESANN}, 2013 (UCI HAR dataset)\\
        {[2]} Menz et al., \textit{Age \& Ageing}, 32(1), 2003 (gait drift)\\
        {[3]} Sculley et al., \textit{NIPS}, 2014 (MLOps, technical debt)\\
        {[4]} \texttt{scikit-learn.org/stable/user\_guide}\\
        {[5]} \texttt{scipy.stats.ks\_2samp} documentation
    \end{columns}
    \normalsize
\end{frame}

\begin{frame}{Questions?}
    \begin{center}
        \Huge Questions?

        \vspace{0.8cm}
        \normalsize
        \textbf{Eduardo Baena}\\
        \texttt{e.baena@northeastern.edu}

        \vspace{0.5cm}
        \textit{EECE5155: Wireless Sensor Networks and the Internet of Things}\\
        Spring 2026

        \vspace{0.4cm}
        \textbf{Real data $\rightarrow$ Feature engineering $\rightarrow$ Train $\rightarrow$ Compare methods $\rightarrow$ Deploy $\rightarrow$ Drift $\rightarrow$ Retrain}

        \vspace{0.3cm}
        {\small \textbf{Take-home due:} one week from today on Canvas}
    \end{center}
\end{frame}

\end{document}
